



\subsection{Descriptive Analysis}\label{sec:descriptives}
We now describe the disadvantaged students targeted by PACE, and their higher education choices absent preferential admissions.\par





\paragraph{Fact 1: Students targeted by PACE score substantially worse on high school standardized tests than regular entrants in selective colleges, and come from poorer households.} Figure \ref{fig:comparison_simce} shows the distribution of standardized tests scores in $10^{th}$ grade among students targeted by PACE and among regular college entrants, standardized in the population of $10^{th}$ graders. Students in targeted schools score 1.42 standard deviations below regular entrants on average. Their median score corresponds to the fifth percentile of scores among regular entrants. Even those who graduate in the top $15\%$ of targeted schools score substantially worse than regular college entrants, $0.89$ standard deviations below on average. Their median score corresponds to the fifteenth percentile of scores among regular entrants. For reference, we draw the average high school standardized test scores in OECD countries: the majority of targeted students score below the OECD mean, the majority of regular entrants score above it.\par 

Table \ref{tab:targetpopulation} shows that students in targeted schools are substantially more disadvantaged than the average Chilean student along several dimensions of socioeconomic status, for example, their family income is half that of the average Chilean student. Family income in this group is $50\%$ of the median household income in Chile, and $32\%$ of the family income of regular entrants, whose average family income of CLP 904,354 per month is $56\%$ above the median Chilean income.









\begin{figure}[ht!]
	\centering
\begin{subfigure}[b]{0.49\textwidth}
\includegraphics[width=\linewidth]{Revision/Figures R/histogram_simce_control_all.png}
 %[scale=0.3]
 \end{subfigure}
 \begin{subfigure}[b]{0.49\textwidth}
\includegraphics[width=\linewidth]{Revision/Figures R/histogram_simce_controltop15_all.png}
 \end{subfigure}
	\vspace{-0.5cm}
	\caption{\label{fig:comparison_simce} Distributions of standardized SIMCE test scores in $10^{th}$ grade. The scores are standardized in the population of $10^{th}$ grade students in 2015. Targeted students are those in schools targeted for PACE, we consider those assigned to the control group. The left panel includes all students in these schools, the right panel includes only those who graduate in the top 15\% of their cohort. Each bar represents 0.20 standard deviations of the distribution of grade 10 test scores in the population of 10 graders. The average score in the OECD is calculated using PISA scores, re-scaled to be comparable to the SIMCE scores (for details see Appendix \ref{sec:PISA}).}
	%\end{minipage}
\end{figure}


 PACE targets a substantially more disadvantaged population than the two most well-known percent plans in the United States. Students in California around the Eligibility in the Local Context preferential admission cutoff have family incomes that are $90\%$ of the median Californian income (\citealp{bleemer2021top}, Table 1) and entrance exam (SAT) scores above the average score among all college applicants (\citealp{bleemer2021top}, Table 1). Of the students targeted by Texas Top Ten, $22-23\%$ are eligible for free or reduced school meals (\citealp{black2023winners}, Table 1), compared with $61\%$ of students in PACE schools who are eligible for welfare programs. The students induced to enroll in a more selective college by the Texas Top Ten have entrance exam scores at the $89^{th}$ statewide percentile. \par 
\paragraph{Fact 2: Absent PACE, only few targeted students attend selective college.} Table \ref{tab:summaryoutcomesbeliefs} describes the educational choices of the typical students targeted by PACE absent PACE. Two thirds of students take the college entrance exam (second row of Table \ref{tab:summaryoutcomesbeliefs}), which aligns nicely with our survey data, where $63\%$ report preparing for it. Even students with very low admission likelihoods prepare for and take the entrance exam (Figure \ref{fig:PSU_admissions_lpoly}). But, as the third row of the table shows, exam scores are well below the national average ($-0.60$ standard deviations). Upon observing their exam scores only $21.0\%$ apply to selective colleges. $11.4\%$ of students are admitted and, in the first year after high school graduation, $8.5\%$ enroll in a selective college , located on average 136 km from their high school. Enrollment in selective colleges is almost equally divided between STEM and non-STEM majors. Students who enroll in selective colleges tend to have college peers who are academically higher-performing than themselves: their average college entrance test score is 0.54 standard deviations above the national average.\par 



\input{Revision/Tables R/summary_stats_outcomes.tex} 


Among students who at baseline are in the top $15\%$ of their school (Panel B in the Table), $86\%$ take the entrance exam; their scores are $0.25$ standard deviations below the average test taker's. Upon observing their score, around half (53\%) of those who took the exam apply to selective colleges. Around a third of students are admitted and around a quarter enroll in a selective college, located on average 128 km from their high school. Top-$15\%$ students have similar rates of enrollment in STEM and non-STEM majors, and the average college entrance test score of college peers is 0.67 standard deviations above the national average.\par 



We observe continuous enrollment in or graduation from a selective college five years since first enrolling, overall and by major type. Panel A of Table \ref{tab:summaryoutcomesbeliefs} shows that $59\%$ of those who enroll in the first year are still continuously enrolled or have graduated after five years. Panel B shows that this figure is slightly larger in the sample of high-performing students ($65\%$). The share of college entrants who persist is higher among those who enroll in non-STEM majors, both in the top-15\% and in the whole sample. \par 



Absent the policy, $26.9\%$ of students in targeted schools enroll in vocational higher education programs, $6.1\%$ in non-selective colleges, and $58.5\%$ do not enroll in higher education. Among the top performing students in targeted schools, $25.4\%$ enroll in vocational higher education programs, $10.6\%$ in non-selective colleges, and $38.4\%$ do not enroll in higher education (Panel B). We report in Appendix Table \ref{tab:summaryoutcomesbeliefsTC} descriptions of outcomes in both treated and control schools.\par 